\section{Introduction}

% Paragraph 1: Introduce DeFi networks, credit exposures, risks, and measurement gaps
Decentralised finance (DeFi) has become a significant and rapidly expanding segment of the crypto‑asset ecosystem, locking billions of dollars in value and promising greater financial inclusion, innovation and efficiency \cite{Adamyk2025}.  DeFi networks comprise a rapidly expanding array of protocols—lending platforms, decentralised exchanges, derivatives, asset‑management products and stablecoins—whose interactions are governed by smart contracts and token flows.  Unlike traditional finance, where credit exposure is created by explicit bilateral contracts, DeFi exposure arises when one protocol holds or locks tokens issued by another; this token‑mediated dependence creates a complex network of interdependencies \cite{DeXposure2025}.  A protocol’s tokens may be adopted as collateral, liquidity or reserves across multiple chains, so the failure of one protocol can propagate contagion to all that hold its tokens \cite{DeXposure2025}.  Recent research using the DeXposure dataset shows that inter‑protocol credit exposure has grown rapidly, is concentrated in a small set of core protocols and exhibits declining network density, highlighting a trend toward systemic concentration \cite{DeXposure2025}.  Yet most empirical studies rely on static snapshots or small sets of lending protocols, so there is a critical need for dynamic, ecosystem‑wide datasets and models to monitor how exposure concentrations shift during market events and how new exposure pathways emerge \cite{DeXposure2025}.

The Financial Stability Board warns that by attempting to replicate traditional financial functions, DeFi inherits and may amplify vulnerabilities such as operational fragilities, liquidity and maturity mismatches, leverage and interconnectedness \cite{FSB2023}.  These vulnerabilities can play out differently because automated liquidation mechanisms and reliance on oracles may trigger rapid deleveraging and fire sales \cite{FSB2023}.  The European Systemic Risk Board notes that, by mid‑2025, links between the crypto sector and the broader financial system had intensified, raising financial‑stability concerns, particularly through spillover channels associated with stablecoins \cite{ESRB2025}.  Academic analyses further underscore that DeFi’s benefits—such as financial inclusion and operational efficiency—are accompanied by smart‑contract vulnerabilities, price volatility, market manipulation and liquidity crises; users have suffered significant losses from flash‑loan attacks and coding errors \cite{Adamyk2025}.  Because assets are locked across multiple protocols, a single failure can cascade through the ecosystem, making early detection of vulnerabilities and robust, real‑time risk management essential \cite{Adamyk2025}.  Despite these mounting risks and the growth of inter‑protocol credit exposures, the systemic implications of DeFi networks remain poorly understood due to data limitations and a lack of standardised risk metrics.  A deeper understanding of these exposures is vital for assessing stability, designing appropriate regulation and developing tools to manage and mitigate contagion in decentralised financial markets.


% Paragraph 2: Introduce the DeXposure dataset and DeXposure‑FM model to address the issue
We evaluate \textbf{DeXposure-FM} on a suite of machine-learning tasks that are both policy-relevant and methodologically challenging, drawing on emerging benchmarks for DeFi network modeling \cite{DeXposure2025}. First, the model forecasts edge-level credit exposures and network-wide risk metrics (such as concentration, connectivity, and density) multiple steps into the future, anticipating how inter-protocol dependencies evolve over time. This task builds on prior approaches in dynamic financial networks and captures key trends observed in DeFi exposure data (e.g., increasing concentration of exposures and declining network density) \cite{DeXposure2025}. Second, we conduct scenario-based stress tests by generating shock-consistent trajectories of losses and contagion under extreme events (such as stablecoin de-pegs or major exchange failures). This allows us to simulate how a sudden shock propagates through the web of token dependencies, echoing recent stress-test analyses of DeFi crises (e.g., the 2023 USDC depeg) \cite{Shah2025}. Third, we assess the model’s ability to impute and denoise missing or noisy time-series data on exposures and TVL, using TVL-weighted error metrics and evaluating the impact on downstream forecasting and link-prediction tasks. This imputation task addresses the ubiquitous data gaps in crypto-asset markets, a challenge where modern deep learning methods (e.g. diffusion-based models) have recently made significant progress by leveraging complex temporal patterns \cite{Wang2025ImputeSurvey}. Collectively, these three tasks extend earlier work ranging from vector autoregressions and low-rank factor models to temporal graph neural networks, providing a rigorous evaluation framework for financial graph models. Our results demonstrate the benefits of a multimodal foundation model approach in capturing complex cross-modal patterns of exposure dynamics, yielding more accurate and robust performance than domain-specific models.

% Paragraph 3: Describe the machine‑learning benchmarks used to evaluate the model
Beyond predictive benchmarks, \textbf{DeXposure-FM} offers financial-economic insights with direct macroprudential relevance. The model yields dynamic measures of protocol-level systemic importance, cross-sector spillover indices, and early-warning indicators that respond to shifts in network concentration and reliance on fragile collateral. Such network-level risk signals can provide advance warning of mounting fragility (for instance, identifying concentrations of risky collateral), in line with recent findings on DeFi lending risks \cite{Bertomeu2024}. These statistics directly address policy calls for AI-driven economic measurement \cite{nber2025ai}, illustrating how a multimodal foundation model trained on rich DeFi data can support macroprudential monitoring, stress testing, and even counterfactual policy analysis. In doing so, DeXposure-FM contributes to the emerging literature on systemic risk in decentralized finance \cite{naifar2025mapping,zbandut2025institutionalizing,Aufiero2025}, highlighting the potential of foundation models as new tools for economic measurement and financial regulation. This synergy between AI and financial stability analysis reflects a broader trend: authorities are beginning to explore AI/ML techniques for systemic risk monitoring and supervision, even as they emphasize the need for model transparency and reliability in high-stakes applications \cite{FSB2024AI}. DeXposure-FM’s design — combining interpretable network metrics with cutting-edge forecasts — aligns with these priorities by providing regulators and researchers an integrated view of evolving risks in DeFi networks.

% Paragraph 4: Explain financial‑economics insights generated by the model
Beyond machine‑learning benchmarks, DeXposure‑FM delivers financial economics insights with direct macroprudential relevance.  The model yields dynamic measures of protocol‑level systemic importance, sector‑to‑sector spillover indices and early‑warning indicators based on shifts in network concentration and dependence on fragile collateral.  These statistics respond to policy calls for **AI‑driven economic measurement**   \cite{nber2025ai} and illustrate how multimodal foundation models trained on DeFi data can support macroprudential monitoring, stress testing and counterfactual policy evaluation while highlighting the conditions under which the resulting measures are robust or warrant caution.  In doing so, DeXposure‑FM contributes to the emerging literature on systemic risk in DeFi  \cite{naifar2025mapping,zbandut2025institutionalizing} and demonstrates the potential of foundation models for economic measurement and financial regulation.

Finally, to foster transparency and further research, we open-source all model weights and code; the training and deployment scripts are available on GitHub: {\url{https://github.com/EVIEHub/DeXposure-FM}}. This open release enables independent validation and extension of our results, addressing concerns that “black-box” AI models in finance can be difficult to scrutinize or trust without access to their inner workings \cite{FSB2024AI}. By making DeXposure-FM publicly available, we hope to facilitate a common platform for DeFi risk analysis and encourage development of AI-driven tools for financial stability monitoring.
